\documentclass[11pt]{article}
\usepackage{icmlko}
\begin{document}

\title{무엇을 다시 뽑는지 정하지 않고 잡음을 쟀다 --- 씨앗을 흔들면 $t{=}59$, 행을 흔들면 $t{=}2.9$}
\author{월드모델 실험실}
\date{노트 607 · 608 · 2026년 8월 5일}
\maketitle

\begin{abstract}
논문 $460$ 이 ``\texttt{media\_push} 는 음수 쪽이 더 넓다''($-0.00462$)고
냈다. 노트 $607$ 이 그것을 정정했다 --- ``그 끝은 \textbf{아이돌 $51$ 행}이
만든 것이고 노트 $591$ 의 \texttt{NOISY\_N}$=200$ 아래다''. \textbf{정정의
근거를 안 쟀다.} 재 보니 \textbf{둘 다 반쯤 틀렸다.}

\begin{center}
\begin{tabular}{lrrr}
\hline
아이돌 짝 & 평균 & \textbf{씨앗} SE & \textbf{표본} SE \\
\hline
\texttt{venue\_prominence} & $+0.280$ & $0.0049$ & $\mathbf{0.0973}$ \\
\texttt{media\_push} & $-0.123$ & $0.0050$ & $\mathbf{0.0880}$ \\
\hline
 & & $t=+59.4$ & $\mathbf{t=+2.9}$ \\
\hline
\end{tabular}
\end{center}

\textbf{둘은 다른 잡음이다.} 씨앗 SE 는 같은 $51$ 행을 붙들어 놓고 모형만
다시 뽑은 것이고, 노트 $591$ 이 말한 잡음은 \emph{행}을 다시 뽑는 쪽이다.
\textbf{표본 SE 가 씨앗 SE 의 $12.6$ 배(중앙)}다.
\end{abstract}

\begin{figure}[t]
\centering
\includegraphics[width=\textwidth]{figs/f1.pdf}
\caption{\textbf{가} --- 두 잡음은 $12.6$ 배 다르다.
\textbf{나} --- 표본 SE 는 $1/\sqrt{n}$ 을 따르지만, $51$ 행짜리도 살아남는다.
\textbf{다} --- 자를 바꾸니 시장 $20$ 건이 $0$ 위로 올라온다.}
\end{figure}

\section{두 잡음}

축-도메인 짝 $42$ 개에서 ``그 축을 가리면 그 도메인이 잃는 $\rho$''를 두
가지로 흔들었다 --- 씨앗 $12$(모형을 다시 뽑는다)와 \textbf{유보 행
부트스트랩}($B{=}2000$ $\times$ 씨앗 $4$).

\begin{center}
\begin{tabular}{lrr}
\hline
 & 씨앗 & \textbf{표본} \\
\hline
유보 크기와의 순위상관 & $-0.729$ & $\mathbf{-0.749}$ \\
$0$ 을 무는 짝 & $2/42$ & $\mathbf{31/42}$ \\
\hline
\end{tabular}
\end{center}

\textbf{둘 다 작은 도메인에서 크지만 크기가 열두 배 다르다.} 씨앗으로만
보면 $42$ 짝 중 $40$ 이 ``확실''이고, 행으로 보면 $11$ 만 남는다.

\begin{quote}
\textbf{씨앗 SE 를 신뢰구간으로 쓰면 열두 배 과신한다.} 씨앗이 답하는
물음은 ``\emph{다시 적합하면} 같은 답이 나오나''(재현성)이고, 행이 답하는
물음은 ``\emph{다른 표본이었어도} 같은 답이 나왔나''(일반화)다. 둘 다
정당한 물음이고, \textbf{틀린 것은 하나를 다른 하나로 읽은 것}이다.
\end{quote}

\section{판정이 갈렸다}

\begin{center}
\begin{tabular}{lrrl}
\hline
짝 & 유보 & 표본 $95\%$ 구간 & \\
\hline
\texttt{venue\_pro}$\cdot$아이돌 & $51$ & $[+0.103, +0.482]$ & \textbf{논문 $460$ 이 옳다} \\
\texttt{media\_push}$\cdot$아이돌 & $51$ & $[-0.303, \mathbf{+0.041}]$ & \textbf{노트 $607$ 이 옳다} \\
\texttt{media\_push}$\cdot$팝업 & $65$ & $[-0.215, \mathbf{+0.077}]$ & 노트 $607$ 이 옳다 \\
\hline
\end{tabular}
\end{center}

\textbf{같은 도메인의 같은 $51$ 행인데 한 축은 살고 한 축은 못 산다.}
논문 $460$ 이 판 환산의 음수 끝으로 쓴 \texttt{media\_push} $-0.1235$ 는
$0$ 을 물고, 노트 $607$ 이 ``같은 이유로 잡음''이라 한
\texttt{venue\_prominence} $+0.2848$ 은 \textbf{$0$ 밖이다.}

\section{행수는 대용물이고, 나쁜 대용물이다}

\begin{center}
\begin{tabular}{lrr}
\hline
 & 구간이 $0$ 밖 & $0$ 을 문다 \\
\hline
유보 $200$ 이상 & $9$ & $\mathbf{15}$ \\
유보 $200$ 미만 & $\mathbf{2}$ & $16$ \\
\hline
\end{tabular}
\end{center}

\textbf{행수로 자르면 위양성 $15$ 개를 통과시키고 진짜 $2$ 개를 버린다.}
$650$ 행짜리 웹툰 \texttt{target\_breadth}($+0.0282$)가 $0$ 을 물고,
$51$ 행짜리 아이돌 \texttt{venue\_prominence}($+0.280$)가 안 문다.
\textbf{행수는 SE 를 통해서만 판정에 닿는다} --- 표본 SE 가
$0.70/\sqrt{n}$ 을 잘 따르지만($\rho{=}-0.749$), \emph{효과 크기}는
행수와 무관하므로 비($t$)는 뒤집힐 수 있다.

\begin{quote}
노트 $607$ 이 만든 \texttt{sideaudit.trim\_small} 은 \textbf{행수로
자른다.} 그 자를 버리고 \textbf{구간으로 자르는} \texttt{ci\_keep} 으로
바꿨다.
\end{quote}

\section{시장 $20$ 건은 해롭지 않다}

\begin{center}
\begin{tabular}{lrr}
\hline
자 & \texttt{entry\_friction} & \texttt{media\_push} \\
\hline
전체 & $[-0.00017, +0.00162]$ & $[\mathbf{-0.00461}, +0.00158]$ \\
행수 $200$ 이상 & $[+0.00015, +0.00162]$ & $[-0.00001, +0.00158]$ \\
\textbf{구간이 $0$ 밖} & $\mathbf{+0.00162}$ & $[+0.00133, +0.00158]$ \\
\hline
\end{tabular}
\end{center}

\textbf{논문 $460$ 이 ``부호도 모른다''고 한 것이 부호가 정해진다.} 두
축 다 $+0.0015$ 언저리이고, 팝업 $7$ 건의 기대치 $+0.0010$(노트 $604$)
\emph{위}다.

\begin{quote}
\textbf{그래도 순서를 뒤집지 않는다.} ① 구간이 $0$ 밖인 짝만 고른 것은
\textbf{유의성 위 선택}이라 남은 값이 위로 부풀고, ② 노트 $606$ 이 미리
적은 한계(가진 데서 뺀 값 $\neq$ 없는 데 넣을 값)가 그대로 살아 있고,
③ 팝업 쪽 $+0.0010$ 은 \emph{다른 자}로 잰 값이다(노트 $584$).
\textbf{바뀐 것은 순서가 아니라 순서의 근거다} --- ``$20$ 건은 해로울 수도
있다''가 사라졌고, 남은 근거는 \textbf{건수가 $7$ 대 $20$ 이고 팝업 $7$
행은 지금 동점이라 채우면 반드시 오른다}는 것뿐이다.
\end{quote}

\section*{조항}

\begin{quote}
\textbf{잡음을 재기 전에 ``무엇을 다시 뽑는가''를 먼저 적는다.} 이 판은
$460$ 개 노트 동안 씨앗 스프레드를 신뢰구간처럼 써 왔다. 그것은
\emph{재현성}의 자이고 \emph{일반화}의 자가 아니다. 도메인 수준에서 둘은
$12.6$ 배 다르다. \textbf{그리고 크기로 자르는 문턱은 그 비를 모르므로,
큰 효과를 억울하게 버리고 작은 효과를 억울하게 통과시킨다.}
\end{quote}

\section*{남은 것}

\begin{itemize}
\item \textbf{판 자체의 구간은 얼마나 좁은가} --- 하네스가 찍는 $95\%$
구간(예: 노트 $596$ 의 $[-0.0097,-0.0075]$)은 \textbf{전부 씨앗}이다.
판은 $3{,}369$ 행이라 표본 SE 가 도메인보다 훨씬 작겠지만 \textbf{비를
안 쟀다.} 이 판의 채택 여럿이 여기 걸려 있다 --- \textbf{다음에 잰다.}
\item \textbf{유의성 위 선택을 안 고쳤다} --- $\S4$ 의 ``구간이 $0$ 밖''
값들은 위로 부푼다. 축소 추정이 필요하다.
\item \textbf{부트스트랩이 행을 독립으로 본다} --- 같은 IP 의 여러 행이
있으면 구간이 실제보다 좁다. 군집 부트스트랩을 안 했다.
\end{itemize}

\end{document}
